\documentclass[a4paper]{article}
\usepackage[fontset=fandol]{ctex}
\usepackage[margin=1.65cm]{geometry}
\usepackage{amsmath,amssymb}
\usepackage{graphicx}
\usepackage{hyperref}
\usepackage{longtable,booktabs}
\usepackage{xcolor}
\usepackage{enumitem}
\usepackage{listings}
\lstset{basicstyle=\ttfamily\small,breaklines=true,columns=fullflexible,keepspaces=true}
\hypersetup{colorlinks=true,linkcolor=blue,urlcolor=blue}
\setlist[itemize]{leftmargin=1.25em,itemsep=0.18em}
\setlength{\parskip}{0.45em}
\setlength{\parindent}{2em}
\sloppy
\title{第 09 讲：Scaling laws 基础}
\author{CS336 Spring 2026 public videos and official course materials}
\date{}
\begin{document}
\maketitle
\tableofcontents
\newpage
\section{本章导读}
本章讨论 \textbf{Scaling laws 基础}。CS336 的第一层目标不是把现成大模型当作黑箱使用，而是把 language model 从输入文本、训练目标、模型结构、数据、系统和评估这些部件重新拆开。只有拆开之后，读者才能判断一个设计选择究竟改变了什么。
设想你只有预算训练一次大模型。Scaling laws（缩放律）的作用，是用一组小实验预测大训练，并在固定 compute 下选择参数量和 token 数，而不是用昂贵试错赌一个配置。
本章对应 CS336 Spring 2026 第 09 个公开视频，并结合课程页提供的可解析资料。视频和课程页给出事实边界；正文把这些材料改写为可自学的教材章节。
阅读本章时，先理解问题为什么存在，再看定义、公式和实现形态，随后通过例题和练习检查自己是否能独立复现这条 reasoning chain。
\begin{figure}[h]
\centering
\includegraphics[width=0.58\linewidth]{source_anchor.jpg}
\caption{第 09 讲的课程图像锚点。}
\end{figure}
\subsection{本章知识路线}
\begin{itemize}
\item 本章首先说明 Scaling laws 基础 为什么是 language modeling pipeline 中不可跳过的一环。
\item 其次定义本章反复使用的术语，并说明这些术语对应的计算对象或系统对象。
\item 随后用公式和伪代码表达主要机制，让概念能够落到可计算的变量上。
\item 最后给出例题、风险讨论和练习，便于读者检查自己是否真正掌握了机制。
\end{itemize}
\subsection{结构图}
下面的图用于帮助读者定位本章的核心结构。先看概念之间的依赖，再回到正文中的公式、伪代码和例题。
\begin{figure}[h]
\centering
\includegraphics[width=0.72\linewidth]{pdf_pages/page-02.jpg}
\caption{官方材料图页：Scaling laws 基础 的课程页/PDF/script 图像锚点。}
\end{figure}
\begin{figure}[h]
\centering

\includegraphics[width=0.72\linewidth]{pdf_pages/page-01.jpg}
\caption{官方材料图页：Scaling laws 基础 的课程页/PDF/script 图像锚点。}
\end{figure}
\section{本章主线}
本章可以沿着三条线索阅读：课程为什么把这个主题放在这里，它改变了哪些变量，以及怎样用小实验检查这些变量。
\subsection{本章要解决的核心问题}
本讲介绍 scaling laws 的基础动机：用小规模实验预测大规模训练，而不是在昂贵大模型上盲目调参。
这一点对应的学习动作是：找出对象，写出变量，说明变量如何被测量。若输入规模、硬件、数据分布或评估协议改变，结论也必须重新检查。
因此，Scaling laws 基础 不应被读成若干名词的集合，而应被读成一组可检验的机制。
\subsection{从机制到资源账本}
核心问题是固定 compute budget 时如何选择参数量 N 和训练 token 数 D。Chinchilla 式分析把模型大小、数据量和 loss 连接到一个可决策的曲面。
这一点对应的学习动作是：找出对象，写出变量，说明变量如何被测量。若输入规模、硬件、数据分布或评估协议改变，结论也必须重新检查。
因此，Scaling laws 基础 不应被读成若干名词的集合，而应被读成一组可检验的机制。
\subsection{从课堂讲解到可复现实验}
scaling law 是经验预测工具，不是自然定律。数据分布变化、优化失败、architecture 改动和评估污染都会让外推失效。
这一点对应的学习动作是：找出对象，写出变量，说明变量如何被测量。若输入规模、硬件、数据分布或评估协议改变，结论也必须重新检查。
因此，Scaling laws 基础 不应被读成若干名词的集合，而应被读成一组可检验的机制。
\section{术语、符号与问题边界}
本章的术语横跨模型、数据和系统。读者需要先知道每个词指向的对象：有些词描述数学目标，有些词描述张量形状，有些词描述硬件瓶颈，还有一些词描述评估协议。术语边界不清，后面的公式和实现就会变成空洞的缩写。
\begin{longtable}{p{0.08\linewidth}p{0.42\linewidth}p{0.42\linewidth}}
\toprule
\textbf{\#} & \textbf{术语} & \textbf{阅读要求}\\
\midrule
1 & scaling law（缩放律） & 需要能说明它指向的对象、参与的公式或代码位置，以及一个典型失败情形。\\
2 & loss（损失） & 需要能说明它指向的对象、参与的公式或代码位置，以及一个典型失败情形。\\
3 & compute-optimal（计算最优） & 需要能说明它指向的对象、参与的公式或代码位置，以及一个典型失败情形。\\
4 & Chinchilla & 需要能说明它指向的对象、参与的公式或代码位置，以及一个典型失败情形。\\
5 & IsoFLOP curve（等 FLOPs 曲线） & 需要能说明它指向的对象、参与的公式或代码位置，以及一个典型失败情形。\\
\bottomrule
\end{longtable}
这些术语之间有依赖关系。例如 tokenization 改变 sequence length，sequence length 又改变 attention FLOPs、KV cache 和 evaluation token budget；parallelism 改变显存占用，也改变通信瓶颈和故障恢复方式。
\subsection{关键词}
下面的关键词在课程材料中反复出现。它们不代替定义，只提示读者本章哪些对象需要在后文持续跟踪。
\begin{itemize}
\item scaling, data, laws, size, batch, performance, param, dataset, small, tokens, over, learning, compute, train
\end{itemize}
\section{为什么 scaling laws 有用}
本节讨论 为什么 scaling laws 有用。本节关心的不是名称本身，而是它把 Scaling laws 基础 中的哪些对象变成可计算、可测量或可优化的量。
\subsection{课程目标与问题边界}
视频把本讲称为离开 systems 后回到 deep learning 的 scaling law 基础。
scaling law 让大规模训练从盲目试错变成可预测实验：用小 runs 估计大 run 的 loss。
真正的用途是决策：给定 compute/data budget，应选多大模型、多少 tokens、什么训练长度。
为什么 scaling laws 有用 的作用是把 Scaling laws 基础 中的一个抽象问题变成可计算对象：哪些变量进入公式，哪些状态进入代码，哪些条件会让结果失效。
\subsection{为什么 scaling laws 有用 的机制}
视频把本讲称为离开 systems 后回到 deep learning 的 scaling law 基础。 这给出本节的对象和边界。读者需要把其中的名词对应到张量、数据记录、训练状态或系统状态，而不是停留在缩写层面。
scaling law 让大规模训练从盲目试错变成可预测实验：用小 runs 估计大 run 的 loss。 因而同一个方法在小规模示例、单机实验和集群训练中的意义并不相同。比较方法时，问题规模、硬件条件、数据分布和评估口径必须同时给出。
真正的用途是决策：给定 compute/data budget，应选多大模型、多少 tokens、什么训练长度。 这使本节具有可检验性：如果一个结论不能在公式、伪代码、profile trace、benchmark 或 ablation 中表现出来，它就还只是直觉。
因此，为什么 scaling laws 有用 应当被理解为一组可以实现和测试的假设。CS336 反复强调 from scratch，正是因为只有能被复现的机制，才可能被稳定改进。
\subsection{资源、效率与效果的关系}
\[L(N,D)=L_\infty + A N^{-\alpha}+B D^{-\beta}\]
符号说明：\(L\) 是 validation loss，\(N\) 是参数量，\(D\) 是训练 tokens，\(\alpha,\beta\) 控制随规模下降的速度。
这个关系式给出了本节最小的数学骨架。读者应检查每个变量的单位、取值范围和可观测性；如果某个量只能间接估计，后续实验就必须说明 proxy 是什么。
例如，validation loss 常被用作模型质量的 proxy，tokens/s 用作吞吐 proxy，Jaccard 或 MinHash 用作近重复 proxy，reward model score 用作偏好 proxy。proxy 不是目标本身；它只在评估协议清楚时才可引用。
\subsection{把课程目标写成检查流程}
本节机制可以写成下面的伪代码。它保留课程材料中的计算顺序和状态变化，但省略了工程代码中的错误处理、并行细节和框架样板。
\begin{lstlisting}[language=Python]
fit_loss_surface(small_runs)
for budget in budgets:
    choose_N_D_that_minimizes_predicted_loss(budget)
\end{lstlisting}
读这段伪代码时，重点不是语法，而是状态如何变化：输入是什么，中间变量是什么，主要计算或数据移动发生在哪里，哪些边界条件会破坏结果。
\subsection{边界条件与常见误解}
\begin{itemize}
\item 缩放律是经验模型，不是自然定律。
\item 分布变化、优化失败、数据质量变化会破坏外推。
\item 不要把小实验中的局部结论自动外推到 frontier-scale；scale、数据分布和硬件拓扑都会改变结论。
\item 不要把 benchmark 或 profile 的单个数字当作机制解释；必须同时记录实验设置、输入形状、版本和评估口径。
\end{itemize}
\subsection{与本章其他概念的关系}
为什么 scaling laws 有用 与本章其他部分的关系是互补的：术语表给出对象名称，公式给出最小数学结构，伪代码给出执行顺序，例题则把这些内容压缩成一个可检查的计算。
\subsection{怎样检验 这一课程目标}
检验 这一课程目标 的第一步是确定测量对象。这里的测量对象通常不是一个抽象名词，而是与 scaling laws、loss、parameters、tokens、predict 有关的张量、样本、请求、奖励、通信操作或评估记录。如果对象没有被具体化，后面的实验就只能得到描述性观察，不能得到可复现结论。
第二步是构造最小输入。最小输入应该足够小，使读者能手算或打印关键中间量；同时又要保留本节机制。对于 这一课程目标，最小输入不追求逼真规模，而追求暴露因果关系：改变一个变量时，哪些输出、成本或错误会随之改变。
第三步是写出资源账本。账本可以是 FLOPs、bytes moved、显存峰值、通信量、token 数、样本数、reward 分布或 benchmark 分数。符号说明：\textbackslash{}(L\textbackslash{}) 是 validation loss，\textbackslash{}(N\textbackslash{}) 是参数量，\textbackslash{}(D\textbackslash{}) 是训练 tokens，\textbackslash{}( lpha,\textbackslash{}beta\textbackslash{}) 控制随规模下降的速度。 这类说明应被转化为单位明确的数量关系；如果数量只能间接观测，就必须说明使用了什么 proxy。
第四步是设置对照。一个好对照只改变一个因素，例如 batch size、sequence length、vocabulary size、attention heads、filter threshold、learning rate、decoding 参数或 verifier。如果多个因素同时改变，实验最多说明系统表现不同，不能说明是哪一个机制在起作用。
第五步是观察失败条件。教材中的成功路径往往最容易记住，但工程中真正决定方法边界的是失败案例。本节至少要记住这些限制：缩放律是经验模型，不是自然定律；分布变化、优化失败、数据质量变化会破坏外推。复习时可以把这些限制改写成反例测试。
最后，把实验结论放回 Scaling laws 基础 的整体结构中。为什么 scaling laws 有用 不是孤立技巧：它会影响训练目标、数据选择、系统成本或评估解释中的至少一项。能说清楚这种影响，才说明读者已经从名词记忆进入机制理解。
\subsection{常见混淆：为什么 scaling laws 有用}
scaling laws（缩放律）用小实验预测大模型 loss 或能力趋势。在“为什么 scaling laws 有用”中，这个概念用于解释 Scaling laws 基础 的一个具体机制。
loss 在本节中应被看作一个技术对象，而不是一个孤立名词。需要判断它指向模型、数据、系统还是评估，并说明它在“为什么 scaling laws 有用”中改变了哪个变量或约束。
parameters 在本节中应被看作一个技术对象，而不是一个孤立名词。需要判断它指向模型、数据、系统还是评估，并说明它在“为什么 scaling laws 有用”中改变了哪个变量或约束。
tokens 在本节中应被看作一个技术对象，而不是一个孤立名词。需要判断它指向模型、数据、系统还是评估，并说明它在“为什么 scaling laws 有用”中改变了哪个变量或约束。
predict 在本节中应被看作一个技术对象，而不是一个孤立名词。需要判断它指向模型、数据、系统还是评估，并说明它在“为什么 scaling laws 有用”中改变了哪个变量或约束。
这些概念的共同点是：只有进入公式、代码或评估协议之后才有明确含义。单独背诵名称不会帮助判断一个训练 run、推理系统或数据 pipeline 是否正确。
\subsection{为什么 scaling laws 有用 的小结}
\begin{itemize}
\item 为什么 scaling laws 有用 的任何性能结论都必须同时报告问题规模和评估口径，否则不可比较。
\item 如果一个方法声称降低主要成本，它通常会把成本转移到另一个变量，例如显存、通信、数据质量、训练稳定性或评估复杂度。
\item 公式和伪代码提供推理骨架；真实系统仍需要 profiler、ablation 和 held-out evaluation。
\end{itemize}
\section{IsoFLOPs 与 compute-optimal tradeoff}
现在进入 IsoFLOPs 与 compute-optimal tradeoff。本节关心的不是名称本身，而是它把 Scaling laws 基础 中的哪些对象变成可计算、可测量或可优化的量。
\subsection{为什么要讨论 IsoFLOPs 与 compute-optimal tradeoff}
IsoFLOP curve 固定训练 compute，在不同 \textbackslash{}(N,D\textbackslash{}) 组合中找最小 loss。
Chinchilla 结论推动了对 undertrained large models 的反思：大模型不一定优于更小但训练更充分的模型。
课程强调这种分析应服务于资源配置，而不是把某条论文曲线当作永恒配方。
IsoFLOPs 与 compute-optimal tradeoff 的作用是把 Scaling laws 基础 中的一个抽象问题变成可计算对象：哪些变量进入公式，哪些状态进入代码，哪些条件会让结果失效。
\subsection{IsoFLOPs 与 compute-optimal tradeoff 的机制}
IsoFLOP curve 固定训练 compute，在不同 \textbackslash{}(N,D\textbackslash{}) 组合中找最小 loss。 这给出本节的对象和边界。读者需要把其中的名词对应到张量、数据记录、训练状态或系统状态，而不是停留在缩写层面。
Chinchilla 结论推动了对 undertrained large models 的反思：大模型不一定优于更小但训练更充分的模型。 因而同一个方法在小规模示例、单机实验和集群训练中的意义并不相同。比较方法时，问题规模、硬件条件、数据分布和评估口径必须同时给出。
课程强调这种分析应服务于资源配置，而不是把某条论文曲线当作永恒配方。 这使本节具有可检验性：如果一个结论不能在公式、伪代码、profile trace、benchmark 或 ablation 中表现出来，它就还只是直觉。
因此，IsoFLOPs 与 compute-optimal tradeoff 应当被理解为一组可以实现和测试的假设。CS336 反复强调 from scratch，正是因为只有能被复现的机制，才可能被稳定改进。
\subsection{IsoFLOPs 与 compute-optimal tradeoff 的数学表达}
\[C\approx 6ND,\qquad (N^\star,D^\star)=\arg\min_{6ND=C}L(N,D)\]
符号说明：\(C\) 是总训练 FLOPs；\(N^\star,D^\star\) 是固定 compute 下预测最优的参数量和 token 数。
这个关系式给出了本节最小的数学骨架。读者应检查每个变量的单位、取值范围和可观测性；如果某个量只能间接估计，后续实验就必须说明 proxy 是什么。
例如，validation loss 常被用作模型质量的 proxy，tokens/s 用作吞吐 proxy，Jaccard 或 MinHash 用作近重复 proxy，reward model score 用作偏好 proxy。proxy 不是目标本身；它只在评估协议清楚时才可引用。
\subsection{从 IsoFLOPs 与 compute-optimal tradeoff 到程序}
本节机制可以写成下面的伪代码。它保留课程材料中的计算顺序和状态变化，但省略了工程代码中的错误处理、并行细节和框架样板。
\begin{lstlisting}[language=Python]
for N in candidate_model_sizes:
    D = compute_budget / (6 * N)
    loss = scaling_law(N, D)
select_min_loss_configuration()
\end{lstlisting}
读这段伪代码时，重点不是语法，而是状态如何变化：输入是什么，中间变量是什么，主要计算或数据移动发生在哪里，哪些边界条件会破坏结果。
\subsection{边界条件与常见误解}
\begin{itemize}
\item 公式里的 6ND 是近似；architecture 和 sequence length 会改变常数。
\item compute-optimal 不等于 latency-optimal 或 serving-cost-optimal。
\item 不要把小实验中的局部结论自动外推到 frontier-scale；scale、数据分布和硬件拓扑都会改变结论。
\item 不要把 benchmark 或 profile 的单个数字当作机制解释；必须同时记录实验设置、输入形状、版本和评估口径。
\end{itemize}
\subsection{与本章其他概念的关系}
IsoFLOPs 与 compute-optimal tradeoff 与本章其他部分的关系是互补的：术语表给出对象名称，公式给出最小数学结构，伪代码给出执行顺序，例题则把这些内容压缩成一个可检查的计算。
\subsection{怎样检验 IsoFLOPs 与 compute-optimal tradeoff}
检验 IsoFLOPs 与 compute-optimal tradeoff 的第一步是确定测量对象。这里的测量对象通常不是一个抽象名词，而是与 IsoFLOPs、compute-optimal、Chinchilla、6ND、budget 有关的张量、样本、请求、奖励、通信操作或评估记录。如果对象没有被具体化，后面的实验就只能得到描述性观察，不能得到可复现结论。
第二步是构造最小输入。最小输入应该足够小，使读者能手算或打印关键中间量；同时又要保留本节机制。对于 IsoFLOPs 与 compute-optimal tradeoff，最小输入不追求逼真规模，而追求暴露因果关系：改变一个变量时，哪些输出、成本或错误会随之改变。
第三步是写出资源账本。账本可以是 FLOPs、bytes moved、显存峰值、通信量、token 数、样本数、reward 分布或 benchmark 分数。符号说明：\textbackslash{}(C\textbackslash{}) 是总训练 FLOPs；\textbackslash{}(N\textasciicircum{}\textbackslash{}star,D\textasciicircum{}\textbackslash{}star\textbackslash{}) 是固定 compute 下预测最优的参数量和 token 数。 这类说明应被转化为单位明确的数量关系；如果数量只能间接观测，就必须说明使用了什么 proxy。
第四步是设置对照。一个好对照只改变一个因素，例如 batch size、sequence length、vocabulary size、attention heads、filter threshold、learning rate、decoding 参数或 verifier。如果多个因素同时改变，实验最多说明系统表现不同，不能说明是哪一个机制在起作用。
第五步是观察失败条件。教材中的成功路径往往最容易记住，但工程中真正决定方法边界的是失败案例。本节至少要记住这些限制：公式里的 6ND 是近似；architecture 和 sequence length 会改变常数；compute-optimal 不等于 latency-optimal 或 serving-cost-optimal。复习时可以把这些限制改写成反例测试。
最后，把实验结论放回 Scaling laws 基础 的整体结构中。IsoFLOPs 与 compute-optimal tradeoff 不是孤立技巧：它会影响训练目标、数据选择、系统成本或评估解释中的至少一项。能说清楚这种影响，才说明读者已经从名词记忆进入机制理解。
\subsection{常见混淆：IsoFLOPs 与 compute-optimal tradeoff}
FLOPs（floating-point operations）是训练和推理成本账本的核心单位。在“IsoFLOPs 与 compute-optimal tradeoff”中，这个概念用于解释 Scaling laws 基础 的一个具体机制。
compute-optimal 在本节中应被看作一个技术对象，而不是一个孤立名词。需要判断它指向模型、数据、系统还是评估，并说明它在“IsoFLOPs 与 compute-optimal tradeoff”中改变了哪个变量或约束。
Chinchilla 风格分析把参数量、token 数和 compute budget 联系起来。在“IsoFLOPs 与 compute-optimal tradeoff”中，这个概念用于解释 Scaling laws 基础 的一个具体机制。
6ND 在本节中应被看作一个技术对象，而不是一个孤立名词。需要判断它指向模型、数据、系统还是评估，并说明它在“IsoFLOPs 与 compute-optimal tradeoff”中改变了哪个变量或约束。
budget 在本节中应被看作一个技术对象，而不是一个孤立名词。需要判断它指向模型、数据、系统还是评估，并说明它在“IsoFLOPs 与 compute-optimal tradeoff”中改变了哪个变量或约束。
这些概念的共同点是：只有进入公式、代码或评估协议之后才有明确含义。单独背诵名称不会帮助判断一个训练 run、推理系统或数据 pipeline 是否正确。
\subsection{IsoFLOPs 与 compute-optimal tradeoff 的小结}
\begin{itemize}
\item IsoFLOPs 与 compute-optimal tradeoff 的任何性能结论都必须同时报告问题规模和评估口径，否则不可比较。
\item 如果一个方法声称降低主要成本，它通常会把成本转移到另一个变量，例如显存、通信、数据质量、训练稳定性或评估复杂度。
\item 公式和伪代码提供推理骨架；真实系统仍需要 profiler、ablation 和 held-out evaluation。
\end{itemize}
\section{从小实验到大训练的风险控制}
现在进入 从小实验到大训练的风险控制。本节关心的不是名称本身，而是它把 Scaling laws 基础 中的哪些对象变成可计算、可测量或可优化的量。
\subsection{为什么要讨论 从小实验到大训练的风险控制}
scaling law 实验要覆盖足够的模型大小、数据量和训练长度，否则拟合不稳定。
validation loss 的测量要保持数据分布一致，并避免污染和重复。
当预测大 run 时，置信区间、残差分析和 sanity checks 比单个点预测更重要。
从小实验到大训练的风险控制 的作用是把 Scaling laws 基础 中的一个抽象问题变成可计算对象：哪些变量进入公式，哪些状态进入代码，哪些条件会让结果失效。
\subsection{从小实验到大训练的风险控制 的机制}
scaling law 实验要覆盖足够的模型大小、数据量和训练长度，否则拟合不稳定。 这给出本节的对象和边界。读者需要把其中的名词对应到张量、数据记录、训练状态或系统状态，而不是停留在缩写层面。
validation loss 的测量要保持数据分布一致，并避免污染和重复。 因而同一个方法在小规模示例、单机实验和集群训练中的意义并不相同。比较方法时，问题规模、硬件条件、数据分布和评估口径必须同时给出。
当预测大 run 时，置信区间、残差分析和 sanity checks 比单个点预测更重要。 这使本节具有可检验性：如果一个结论不能在公式、伪代码、profile trace、benchmark 或 ablation 中表现出来，它就还只是直觉。
因此，从小实验到大训练的风险控制 应当被理解为一组可以实现和测试的假设。CS336 反复强调 from scratch，正是因为只有能被复现的机制，才可能被稳定改进。
\subsection{从小实验到大训练的风险控制 的数学表达}
\[\hat L_{\text{large}}\pm z\sigma_{\text{residual}}\]
符号说明：\(\hat L_{\text{large}}\) 是预测 loss，\(\sigma_{\text{residual}}\) 是小实验拟合残差的尺度。
这个关系式给出了本节最小的数学骨架。读者应检查每个变量的单位、取值范围和可观测性；如果某个量只能间接估计，后续实验就必须说明 proxy 是什么。
例如，validation loss 常被用作模型质量的 proxy，tokens/s 用作吞吐 proxy，Jaccard 或 MinHash 用作近重复 proxy，reward model score 用作偏好 proxy。proxy 不是目标本身；它只在评估协议清楚时才可引用。
\subsection{从 从小实验到大训练的风险控制 到程序}
本节机制可以写成下面的伪代码。它保留课程材料中的计算顺序和状态变化，但省略了工程代码中的错误处理、并行细节和框架样板。
\begin{lstlisting}[language=Python]
runs = launch_grid(model_sizes, token_counts)
fit = robust_fit(runs)
plot_residuals_and_refuse_bad_extrapolation(fit)
\end{lstlisting}
读这段伪代码时，重点不是语法，而是状态如何变化：输入是什么，中间变量是什么，主要计算或数据移动发生在哪里，哪些边界条件会破坏结果。
\subsection{边界条件与常见误解}
\begin{itemize}
\item 过度外推会把小模型 regime 的现象误用于大模型。
\item 数据 pipeline 改变后，旧 scaling law 需要重新校准。
\item 不要把小实验中的局部结论自动外推到 frontier-scale；scale、数据分布和硬件拓扑都会改变结论。
\item 不要把 benchmark 或 profile 的单个数字当作机制解释；必须同时记录实验设置、输入形状、版本和评估口径。
\end{itemize}
\subsection{与本章其他概念的关系}
从小实验到大训练的风险控制 与本章其他部分的关系是互补的：术语表给出对象名称，公式给出最小数学结构，伪代码给出执行顺序，例题则把这些内容压缩成一个可检查的计算。
\subsection{怎样检验 从小实验到大训练的风险控制}
检验 从小实验到大训练的风险控制 的第一步是确定测量对象。这里的测量对象通常不是一个抽象名词，而是与 extrapolation、residual、validation、grid、forecast 有关的张量、样本、请求、奖励、通信操作或评估记录。如果对象没有被具体化，后面的实验就只能得到描述性观察，不能得到可复现结论。
第二步是构造最小输入。最小输入应该足够小，使读者能手算或打印关键中间量；同时又要保留本节机制。对于 从小实验到大训练的风险控制，最小输入不追求逼真规模，而追求暴露因果关系：改变一个变量时，哪些输出、成本或错误会随之改变。
第三步是写出资源账本。账本可以是 FLOPs、bytes moved、显存峰值、通信量、token 数、样本数、reward 分布或 benchmark 分数。符号说明：\textbackslash{}(\textbackslash{}hat L\_\{ ext\{large\}\}\textbackslash{}) 是预测 loss，\textbackslash{}(\textbackslash{}sigma\_\{ ext\{residual\}\}\textbackslash{}) 是小实验拟合残差的尺度。 这类说明应被转化为单位明确的数量关系；如果数量只能间接观测，就必须说明使用了什么 proxy。
第四步是设置对照。一个好对照只改变一个因素，例如 batch size、sequence length、vocabulary size、attention heads、filter threshold、learning rate、decoding 参数或 verifier。如果多个因素同时改变，实验最多说明系统表现不同，不能说明是哪一个机制在起作用。
第五步是观察失败条件。教材中的成功路径往往最容易记住，但工程中真正决定方法边界的是失败案例。本节至少要记住这些限制：过度外推会把小模型 regime 的现象误用于大模型；数据 pipeline 改变后，旧 scaling law 需要重新校准。复习时可以把这些限制改写成反例测试。
最后，把实验结论放回 Scaling laws 基础 的整体结构中。从小实验到大训练的风险控制 不是孤立技巧：它会影响训练目标、数据选择、系统成本或评估解释中的至少一项。能说清楚这种影响，才说明读者已经从名词记忆进入机制理解。
\subsection{常见混淆：从小实验到大训练的风险控制}
extrapolation 在本节中应被看作一个技术对象，而不是一个孤立名词。需要判断它指向模型、数据、系统还是评估，并说明它在“从小实验到大训练的风险控制”中改变了哪个变量或约束。
residual 在本节中应被看作一个技术对象，而不是一个孤立名词。需要判断它指向模型、数据、系统还是评估，并说明它在“从小实验到大训练的风险控制”中改变了哪个变量或约束。
validation 在本节中应被看作一个技术对象，而不是一个孤立名词。需要判断它指向模型、数据、系统还是评估，并说明它在“从小实验到大训练的风险控制”中改变了哪个变量或约束。
grid 在本节中应被看作一个技术对象，而不是一个孤立名词。需要判断它指向模型、数据、系统还是评估，并说明它在“从小实验到大训练的风险控制”中改变了哪个变量或约束。
forecast 在本节中应被看作一个技术对象，而不是一个孤立名词。需要判断它指向模型、数据、系统还是评估，并说明它在“从小实验到大训练的风险控制”中改变了哪个变量或约束。
这些概念的共同点是：只有进入公式、代码或评估协议之后才有明确含义。单独背诵名称不会帮助判断一个训练 run、推理系统或数据 pipeline 是否正确。
\subsection{从小实验到大训练的风险控制 的小结}
\begin{itemize}
\item 从小实验到大训练的风险控制 的任何性能结论都必须同时报告问题规模和评估口径，否则不可比较。
\item 如果一个方法声称降低主要成本，它通常会把成本转移到另一个变量，例如显存、通信、数据质量、训练稳定性或评估复杂度。
\item 公式和伪代码提供推理骨架；真实系统仍需要 profiler、ablation 和 held-out evaluation。
\end{itemize}
\section{综合讨论：从局部机制到完整系统}
本章的主题是 Scaling laws 基础，但它不应被读成几个相邻知识点的拼接。为什么 scaling laws 有用、IsoFLOPs 与 compute-optimal tradeoff、从小实验到大训练的风险控制 共同回答同一个问题：语言模型系统中哪些对象可以被定义，哪些成本可以被估算，哪些结论必须通过实验或评估来确认。
这些对象包括 scaling law（缩放律）、loss（损失）、compute-optimal（计算最优）、Chinchilla、IsoFLOP curve（等 FLOPs 曲线）。读者可以把它们看成一组接口：有的接口连接文本和 token，有的连接模型结构和张量计算，有的连接数据样本和训练目标，有的连接离线评估和在线服务。接口一旦改变，后面的成本、错误类型和优化空间也会随之改变。
\subsection{为什么 scaling laws 有用 在系统中的位置}
为什么 scaling laws 有用 首先提供一个局部解释：视频把本讲称为离开 systems 后回到 deep learning 的 scaling law 基础。 这类解释的价值在于把直觉压缩成可操作对象。一旦对象明确，读者就可以追问它的输入、输出、中间状态和失败条件，而不是只记住方法名。
其次，为什么 scaling laws 有用 会改变至少一个资源或评估变量。它可能改变 sequence length、activation size、memory bandwidth、communication volume、data mixture、reward distribution 或 benchmark protocol。这些变量在小例子里可以手算，在真实训练和服务系统里则需要 profiling、logging 和 held-out evaluation。
最后，为什么 scaling laws 有用 不能脱离边界条件理解。一个关键 caveat 是：缩放律是经验模型，不是自然定律。。把 caveat 写出来不是为了削弱结论，而是为了说明结论在哪些输入、硬件、数据和评估条件下才成立。
\subsection{IsoFLOPs 与 compute-optimal tradeoff 在系统中的位置}
IsoFLOPs 与 compute-optimal tradeoff 首先提供一个局部解释：IsoFLOP curve 固定训练 compute，在不同 \textbackslash{}(N,D\textbackslash{}) 组合中找最小 loss。 这类解释的价值在于把直觉压缩成可操作对象。一旦对象明确，读者就可以追问它的输入、输出、中间状态和失败条件，而不是只记住方法名。
其次，IsoFLOPs 与 compute-optimal tradeoff 会改变至少一个资源或评估变量。它可能改变 sequence length、activation size、memory bandwidth、communication volume、data mixture、reward distribution 或 benchmark protocol。这些变量在小例子里可以手算，在真实训练和服务系统里则需要 profiling、logging 和 held-out evaluation。
最后，IsoFLOPs 与 compute-optimal tradeoff 不能脱离边界条件理解。一个关键 caveat 是：公式里的 6ND 是近似；architecture 和 sequence length 会改变常数。。把 caveat 写出来不是为了削弱结论，而是为了说明结论在哪些输入、硬件、数据和评估条件下才成立。
\subsection{从小实验到大训练的风险控制 在系统中的位置}
从小实验到大训练的风险控制 首先提供一个局部解释：scaling law 实验要覆盖足够的模型大小、数据量和训练长度，否则拟合不稳定。 这类解释的价值在于把直觉压缩成可操作对象。一旦对象明确，读者就可以追问它的输入、输出、中间状态和失败条件，而不是只记住方法名。
其次，从小实验到大训练的风险控制 会改变至少一个资源或评估变量。它可能改变 sequence length、activation size、memory bandwidth、communication volume、data mixture、reward distribution 或 benchmark protocol。这些变量在小例子里可以手算，在真实训练和服务系统里则需要 profiling、logging 和 held-out evaluation。
最后，从小实验到大训练的风险控制 不能脱离边界条件理解。一个关键 caveat 是：过度外推会把小模型 regime 的现象误用于大模型。。把 caveat 写出来不是为了削弱结论，而是为了说明结论在哪些输入、硬件、数据和评估条件下才成立。
\subsection{把本章知识用于读论文和复现实验}
读相关论文时，可以把本章变成一个检查表。第一，论文是否清楚定义了输入规模、模型规模、数据来源和评估协议；第二，论文声称的改进落在哪个变量上；第三，作者是否报告了会让结论失效的设置；第四，实验是否能分离模型结构、数据质量、优化设置和系统实现的影响。
复现实验时，也应避免直接追求大规模。先用最小程序复现一个公式、一个伪代码分支或一个 profile 现象，再扩大输入规模。这样做的原因很简单：如果小规模程序无法解释中间量，大规模训练只会把错误隐藏在日志和随机波动里。
把这些检查合在一起，本章给出的不是一个固定答案，而是一种阅读 CS336 的方法：每个模型机制都要落到变量，每个变量都要落到测量，每个测量都要说明协议，每个协议都要承认边界。
因此，本章中的任何一句结论都可以被改写成一个可引用的完整句式：在给定数据、模型、硬件和评估协议下，某个机制改变了某个变量，并通过某个观测指标表现出来。这样的句子比单独背诵术语更长，但它包含了自学、复现和引用时真正需要的信息。
把本章用于自学时，可以按三遍阅读。第一遍只追踪对象和定义，确认每个术语到底指向文本、token、张量、样本、请求还是评估记录。第二遍追踪公式和伪代码，确认每一步改变了哪些状态。第三遍追踪实验和 caveat，确认哪些结论只在特定规模、数据或硬件条件下成立。
把本章用于复习时，则应反过来操作：先从一个失败案例或异常指标出发，再回到相应机制。loss 曲线异常可能指向数据、优化或模型容量；吞吐异常可能指向 kernel、通信或调度；benchmark 异常可能指向 contamination、prompt protocol 或采样设置。能够从异常反推机制，是掌握本章的实用标准。
\section{例题：把概念变成可计算检查}
\subsection{固定 compute 下选 N 与 D}
给定 `C≈6ND`，如果模型参数 `N` 翻倍，而 compute 固定，则训练 tokens `D` 约减半。
\begin{itemize}
\item 大模型 token 不足会 undertrained。
\item 小模型 token 充足但容量可能限制 loss。
\item IsoFLOP 实验通过多个 `N,D` 点拟合 loss 曲面。
\end{itemize}
\textbf{结论。} compute-optimal 是一个受数据、模型和目标指标共同影响的选择问题。
这个例题的目的不是替代完整工程实现，而是给出一种最小推理形式：把问题转成变量，写出近似公式或伪代码，再说明近似何时失效。
\section{实验设计：从小例子到可复现结论}
学习 Scaling laws 基础 不能停在概念层。一个教材化结论至少需要四件事：可构造的输入、可观测的中间量、可比较的指标，以及清楚的失败条件。下面的实验设计不是要求训练大模型，而是要求读者用小规模程序复现本章机制。
\subsection{最小输入}
先构造只触发本章核心机制的小输入。例如 tokenizer 章节可以用几个包含 rare words 和 Unicode 字符的字符串；GPU 章节可以用一个 elementwise kernel 和一个 GEMM；data 章节可以用十几篇近重复文档。小输入的价值在于它让错误可见。
\subsection{资源账本}
对同一输入写下 FLOPs、bytes moved、显存峰值、通信量或样本数量的估算。估算不需要一开始就精确，但必须说明单位和近似假设。随后用 profiler、benchmark、validation loss 或人工检查去验证估算是否同量级。
\subsection{单变量消融}
一次只改变一个因素，例如 vocabulary size、head 数、batch size、filter threshold、reward scale、decoding temperature 或 GPU 数量。若多个变量同时变化，实验只能说明系统变了，不能说明机制是什么。
\subsection{失败条件}
最后要刻意触发失败：极端 context length、错误 dtype、过小 batch、污染 benchmark、低质量数据或不可靠 verifier。失败条件常常比成功曲线更能说明一个方法的边界。
\subsection{记录与报告}
实验记录应能让另一个读者复现同一结论。最少要写清楚输入构造、代码版本、随机种子、硬件、batch shape、数据版本、评价指标和异常处理方式。报告结论时不要只写“更快”或“更好”，而要说明快在哪里、好在哪个指标上，以及是否牺牲了内存、稳定性、数据质量或可解释性。
常见报告错误是把局部现象写成普遍规律。例如一次小模型实验里的 loss 改善，不等于更大模型或不同数据分布上也会改善；一次 kernel benchmark 的吞吐提升，不等于端到端 serving latency 一定降低。教材中的实验报告应始终把结论限定在可复现条件内。
\section{课程资料与回看路径}
教材正文已经把主要概念、公式和实现逻辑组织成连续叙述；本节只提供回到课程材料的阅读路径。读者复习时可以先完成前文练习，再按下面的时间段或资料组核对细节。
\subsection{视频回看路径}
\begin{longtable}{p{0.20\linewidth}p{0.46\linewidth}p{0.26\linewidth}}
\toprule
\textbf{时间} & \textbf{关键词} & \textbf{复习任务}\\
\midrule
00:00:06-00:06:49 & scaling, very, going, that's, laws, scale, your & 回看定义、例子、推导或 caveat 在该段中的位置。\\
00:20:24-00:26:39 & data, scaling, than, small, that's, mixture, laws & 回看定义、例子、推导或 caveat 在该段中的位置。\\
00:39:11-00:45:44 & batch, number, going, size, your, parameters, point & 回看定义、例子、推导或 caveat 在该段中的位置。\\
00:52:07-00:58:35 & very, going, think, scaling, data, it's, your & 回看定义、例子、推导或 caveat 在该段中的位置。\\
01:11:20-01:17:53 & think, scaling, chinchilla, compute, very, training, scale & 回看定义、例子、推导或 caveat 在该段中的位置。\\
\bottomrule
\end{longtable}
这些时间段用于定位视频中的讲解顺序。严肃引用时，应回到视频片段确认讲者表述、板书或幻灯片内容。
回看视频时，不必逐字抄录字幕。更有效的方法是把每个时间段判定为定义、例子、推导、代码解释或 caveat，再把它放回本章对应小节。
\subsection{课程资料阅读路径}
\begin{longtable}{p{0.14\linewidth}p{0.48\linewidth}p{0.30\linewidth}}
\toprule
\textbf{资料组} & \textbf{标题/页块} & \textbf{关键词}\\
\midrule
official\_01 & Lecture 9 S C A LIN G LAW S - B A SIC S CS336 & scaling, laws, your, together, simple, predictive, tune\\
official\_02 & Earliest (data) scaling law paper – 1993 & scaling, data, early, earliest, history, laws, hestness\\
official\_03 & Hestness II Very ahead of its time.. “Emergence” Scaling by compute Speed = accuracy & scaling, data, laws, performance, dataset, size, error\\
official\_05 & Side note – ‘Value of parameters’ With MoEs, we expect parameter scaling to change & scaling, batch, size, critical, data, joint, over\\
official\_06 & Explanation 2 Non-embedding vs total param choice + small nonlinearities & param, tokens, data, scaling, compute, chinchilla, small\\
\bottomrule
\end{longtable}
课程资料通常给出更稳定的标题、公式、图或代码结构；视频则补充动机和口头解释。两者合起来构成本章的事实边界。
阅读资料时，可以把每个资料组拆成三个对象：一个定义，一个公式或伪代码动作，一个边界条件。这样比逐字翻译页面或脚本更容易形成可用知识。
\section{实践说明、历史位置与风险}
Scaling laws 基础 在 CS336 的课程结构中连接基础机制和规模化问题。基础机制解释方法为什么成立；规模化问题检验它在更大模型、更长上下文、更复杂数据或真实 serving workload 下是否仍成立。
历史位置不等于模型年表。更重要的是识别问题如何迁移：早期方法解决小规模建模问题，现代方法常常解决同一问题在大规模数据、GPU 集群、长上下文或人类偏好约束下的新形态。
\begin{itemize}
\item 资料版本限制（version risk）：课程页、公开视频和配套资料可能在学期中更新；引用时应注明访问日期或资料版本。
\item 测量口径限制（measurement risk）：FLOPs、tokens/s、benchmark score、reward 等数字只有在协议完整时可比较。
\item 规模外推限制（scaling risk）：小模型或小 batch 的机制不必然外推到 frontier-scale。
\item 数据分布限制（data risk）：数据来源、过滤和去重策略会改变模型行为，也会改变 evaluation contamination 风险。
\end{itemize}
\section{复习要点}
下面的问题把本章内容压缩为可反复检查的条目。每个条目都应能回到前文的解释、公式、伪代码或例题。
\begin{longtable}{p{0.16\linewidth}p{0.38\linewidth}p{0.38\linewidth}}
\toprule
\textbf{条目} & \textbf{对象} & \textbf{检查动作}\\
\midrule
条目 1 & scaling law（缩放律） & 定义它；写出一个最小例子；说明一个 failure mode；指出它影响哪个资源或评估账本。\\
条目 2 & loss（损失） & 定义它；写出一个最小例子；说明一个 failure mode；指出它影响哪个资源或评估账本。\\
条目 3 & compute-optimal（计算最优） & 定义它；写出一个最小例子；说明一个 failure mode；指出它影响哪个资源或评估账本。\\
条目 4 & Chinchilla & 定义它；写出一个最小例子；说明一个 failure mode；指出它影响哪个资源或评估账本。\\
条目 5 & IsoFLOP curve（等 FLOPs 曲线） & 定义它；写出一个最小例子；说明一个 failure mode；指出它影响哪个资源或评估账本。\\
\bottomrule
\end{longtable}
\subsection{综合自测}
\begin{itemize}
\item 我能否不用课程原句，独立解释本章中心问题？
\item 我能否把本章至少一个公式改写成可运行的估算函数？
\item 我能否指出一个看似改进但实际转移成本的方法？
\item 我能否回到视频和课程资料，找到支持本章关键论断的位置？
\item 我能否设计一个最小 ablation，验证本章一个 caveat？
\end{itemize}
\section{总结与延伸}
\subsection{本章小结}
\begin{itemize}
\item Scaling laws 基础 的核心不是记忆术语，而是理解对象、变量和机制之间的关系。
\item 视频给出讲解顺序，课程资料提供可核对的公式、图、代码或页面结构。
\item 复习时应把每个术语连接到至少一个变量、一个实现动作和一个失败模式。
\end{itemize}
\subsection{延伸解释}
本章中的延伸解释用于连接课程材料中的概念，例如说明公式里的 proxy、解释系统 tradeoff，或把多个材料片段串成一个完整推理链。若需要严肃引用，应回到课程视频和配套资料核对原始表述。
\subsection{练习}
\begin{itemize}
\item 定义题：用中英双语解释本章任意五个重要术语，并说明它们属于 compute、memory、data、optimization、evaluation 还是 deployment 账本。
\item 推导题：选择本章一个公式，逐个解释符号、单位和近似假设，并说明哪个变量最难在真实系统中测量。
\item 实现题：把本章伪代码改写成一个最小 Python sanity check，不要求训练大模型，但要输出可检查的中间量。
\item 资料题：从“课程资料与回看路径”中选择一个视频时间段，回到原始视频核对讲者例子，并记录是否存在字幕误差。
\item 评估题：为本章方法设计一个 ablation 或 benchmark，说明控制变量、数据版本、指标和 failure criteria。
\item 边界题：说明本章一个结论在更长 context、更大 batch、更多 GPU 或更复杂数据分布下可能如何失效。
\item 综合题：把本章内容和前一章或后一章连接起来，写出一个端到端训练/推理 pipeline 中的依赖关系。
\item 批判题：指出一个容易被营销材料夸大的说法，并用本章的资源账本或评估协议反驳它。
\end{itemize}
\end{document}
